
\section{非线性\BookRecordIndex{T0209}柯西--黎曼方程}\label{sec:2-2}
\zhindex{feixianxingkexi-limanfangcheng}{非线性柯西--黎曼方程}

设 $(M,J)$ 是殆复流形，$(\Sigma,j)$ 是黎曼曲面。若光滑映射 $u:\Sigma\to M$ 的微分
$du$ 关于 $j$ 和 $J$ 是复线性的，即
\[
  J\circ du=du\circ j,
\]
就称 $u$ 是一条 \BookRecordIndex{T0003}$J$-全纯曲线，简称 $J$-曲线。有时把这个方程写成
$\bar\partial_J(u)=0$ 更方便，其中 $1$-形式
\[
  \bar\partial_J(u):=\frac12\bigl(du+J\circ du\circ j\bigr)
  \in\Omega^{0,1}(\Sigma,u^*TM)
\]
是 $du$ 关于\BookRecordIndex{T0088}殆复结构 $J$ 的复反线性部分。它取值于拉回切丛
\[
  u^*TM:=\{(z,v)\mid z\in\Sigma,\ v\in T_{u(z)}M\}.
\]
用更抽象的语言说，在从 $\Sigma$ 到 $M$ 的光滑映射空间
$\mathcal B:=C^\infty(\Sigma,M)$ 上，有一个无限维向量丛
$\mathcal E\to\mathcal B$；它在 $u\in\mathcal B$ 上的纤维是
$\mathcal E_u:=\Omega^{0,1}(\Sigma,u^*TM)$。非线性算子
\[
  \mathcal B\longrightarrow\mathcal E,
  \qquad u\longmapsto\bigl(u,\bar\partial_J(u)\bigr)
\]
是该丛的一个截面，其零点集就是 \BookRecordIndex{T0003}$J$-全纯曲线空间。当 $\Sigma$ 的亏格为 $0$ 或 $1$ 时，
值得注意的是，在 $\Sigma$ 的复自同构群 $G$ 对 $\mathcal B$ 和 $\mathcal E$ 的作用下，
投影 $\mathcal E\to\mathcal B$ 和截面 $\bar\partial_J$ 都是等变的。下面会详细讨论这一点。
第 3 章尤其将说明，在对 $\mathcal B$ 和 $\mathcal E$ 取合适的索伯列夫完备化后，
$\bar\partial_J$ 是一个弗雷德霍姆截面。方程的整体形式有助于研究所有 \BookRecordIndex{T0003}$J$-全纯曲线的
\BookRecordIndex{T0246}模空间。本章主要考察 \BookRecordIndex{T0003}$J$-全纯曲线的局部性质，首先在局部坐标中写出方程。

\paragraph{局部坐标。}
可积性定理蕴含 $\Sigma$ 上存在由全纯坐标卡
$\phi_\alpha:U_\alpha\to\mathbb C$ 组成的图册
$\{U_\alpha,\phi_\alpha\}_\alpha$。这表示
\[
  d\phi_\alpha\circ j=i\,d\phi_\alpha,
\]
因而转移映射 $\phi_\beta\circ\phi_\alpha^{-1}$ 全纯。这类局部坐标也称为\BookRecordIndex{T0149}共形坐标。
回忆，$\mathbb C$ 的开子集之间的映射 $\phi$ 是共形的，即保持角度和定向，当且仅当它
全纯。光滑映射 $u:\Sigma\to M$ 是 $J$-全纯的，当且仅当它的局部坐标表示
\[
  u_\alpha=u\circ\phi_\alpha^{-1}:\phi_\alpha(U_\alpha)\to M
\]
关于开集 $\phi_\alpha(U_\alpha)\subset\mathbb C$ 上的标准复结构 $i$ 是 $J$-全纯的。

在 $\Sigma$ 上的\BookRecordIndex{T0149}共形坐标 $z=s+it$ 中，$1$-形式 $\bar\partial_J(u)$ 写成
\begin{align*}
  \bar\partial_J(u_\alpha)
  ={}&\frac12\bigl(\partial_su_\alpha+J(u_\alpha)\partial_tu_\alpha\bigr)\,ds\\
     &+\frac12\bigl(\partial_tu_\alpha-J(u_\alpha)\partial_su_\alpha\bigr)\,dt.
\end{align*}
因此，$u$ 是 \BookRecordIndex{T0003}$J$-全纯曲线，当且仅当它在\BookRecordIndex{T0149}共形坐标中满足非线性一阶偏微分方程
\begin{equation}\label{eq:2-2-1}
  \partial_su_\alpha+J(u_\alpha)\partial_tu_\alpha=0.
\end{equation}
这里 $\partial_s$ 表示偏导数 $\partial/\partial s$。若同一化
$\mathbb C^n=\mathbb R^{2n}$，则标准复结构 $J=i=J_0$ 是常矩阵
\[
  J_0=
  \begin{pmatrix}
    0&-\mathbf1\\
    \mathbf1&0
  \end{pmatrix}.
\]
此时，对光滑映射 $u=f+ig:\mathbb C\to\mathbb C^n$，方程
\eqref{eq:2-2-1} 等价于熟悉形式的\BookRecordIndex{T0209}柯西--黎曼方程
\[
  \partial_sf=\partial_tg,
  \qquad
  \partial_sg=-\partial_tf.
\]
所以 $J_0$-全纯曲线就是通常意义下的全纯曲线。

\paragraph{\BookRecordIndex{T0249}能量恒等式。}
设 $(M,\omega)$ 是\BookRecordIndex{T0089}殆辛流形（即 $\omega$ 非退化但不一定闭），
$J\in\mathcal J_\tau(M)$ 是被 $\omega$ \BookRecordIndex{T0372}驯服的\BookRecordIndex{T0088}殆复结构，并且
$(\Sigma,j_\Sigma,d\operatorname{vol}_\Sigma)$ 是紧黎曼曲面。由
\eqref{eq:2-1-1}，$\omega$ 和 $J$ 确定 $M$ 上的度量 $g_J$；$j_\Sigma$ 和
$d\operatorname{vol}_\Sigma$ 则确定黎曼曲面上的度量。每条 \BookRecordIndex{T0003}$J$-全纯曲线
$u:\Sigma\to M$ 关于 $g_J$ 都是共形的，即其微分保持角度；等价地，它把内积保持到
一个公共的正因子。当 $M$ 为 $2$ 维时，逆命题也成立。

光滑映射 $u:\Sigma\to M$ 的能量定义为 $1$-形式
$du\in\Omega^1(\Sigma,u^*TM)$ 的 $L^2$-范数平方：
\[
  E(u):=\frac12\int_\Sigma |du|_J^2\,d\operatorname{vol}_\Sigma.
\]
这里实线性映射 $L:=du(z):T_z\Sigma\to T_{u(z)}M$ 的范数定义为
\[
  |L|_J:=|\zeta|^{-1}
  \sqrt{|L(\zeta)|_J^2+|L(j_\Sigma\zeta)|_J^2},
  \qquad 0\ne\zeta\in T_z\Sigma.
\]
由下面的习题 2.2.3，右边与 $\zeta$ 的选择无关。$\Sigma$ 上\BookRecordIndex{T0149}共形坐标中被积函数的显式
公式见下面引理 2.2.1 的证明。

必须注意：虽然\BookRecordIndex{T0251}能量密度 $|du|_J^2$ 依赖于 $\Sigma$ 上的度量，能量 $E(u)$ 却只依赖
复结构 $j_\Sigma$，而不依赖体积形式。要看出这一点，注意对任意埃尔米特向量丛
$E\to\Sigma$ 和任意 $1$-形式 $\alpha\in\Omega^1(\Sigma,E)$，都有
\[
  |\alpha|^2d\operatorname{vol}_\Sigma
  =\langle\alpha\wedge *\alpha\rangle
  =-\langle\alpha\wedge(\alpha\circ j_\Sigma)\rangle.
\]
对一般光滑映射 $u:\Sigma\to M$，能量 $E(u)$ 还依赖于 $M$ 上的度量 $g_J$。不过，
下面将说明，在\BookRecordIndex{T0358}辛流形中的 \BookRecordIndex{T0003}$J$-全纯曲线这一情形，能量是一个拓扑不变量，只依赖曲线模其
边界的同调类。

\begin{lemma}\label{lem:2-2-1}
设 $\omega$ 是光滑流形 $M$ 上的非退化 $2$-形式。若 $J$ 被 $\omega$ \BookRecordIndex{T0372}驯服，则每条
\BookRecordIndex{T0003}$J$-全纯曲线 $u:\Sigma\to M$ 都满足\BookRecordIndex{T0249}能量恒等式
\begin{equation}\label{eq:2-2-2}
  E(u)=\int_\Sigma u^*\omega.
\end{equation}
若 $J$ 与 $\omega$ \BookRecordIndex{T0348}相容，则每个光滑映射 $u:\Sigma\to M$ 都满足
\begin{equation}\label{eq:2-2-3}
  E(u)=\int_\Sigma|\bar\partial_J(u)|_J^2\,
  d\operatorname{vol}_\Sigma+\int_\Sigma u^*\omega.
\end{equation}
\end{lemma}

\begin{proof}
选择\BookRecordIndex{T0149}共形坐标 $z=s+it$ 后，不失一般性可设 $\Sigma$ 是 $\mathbb C$ 的开子集。此时
\begin{align*}
  \frac12|du|_J^2d\operatorname{vol}_\Sigma
  &={\frac12}\bigl(|\partial_su|_J^2+|\partial_tu|_J^2\bigr)\,ds\wedge dt\\
  &={\frac12}|\partial_su+J\partial_tu|_J^2\,ds\wedge dt
    -\langle\partial_su,J\partial_tu\rangle_J\,ds\wedge dt\\
  &=|\bar\partial_J(u)|_J^2d\operatorname{vol}_\Sigma\\
  &\quad+{\frac12}\bigl(\omega(\partial_su,\partial_tu)
    +\omega(J\partial_su,J\partial_tu)\bigr)\,ds\wedge dt.
\end{align*}
在 $\omega$-\BookRecordIndex{T0348}相容情形，右边最后一项等于 $u^*\omega$。若 $J$ 被 $\omega$ \BookRecordIndex{T0372}驯服且
$\partial_su+J\partial_tu=0$，则右边第一项消失，最后一项也等于 $u^*\omega$。
\end{proof}

引理 2.2.1 表明，若 $\omega$ 闭，$J$ 与 $\omega$ \BookRecordIndex{T0348}相容且 $\Sigma$ 闭，则
\BookRecordIndex{T0003}$J$-全纯曲线在其同调类中使能量最小，因而是\BookRecordIndex{T0105}调和映射。如果 $J$ 只是被 $\omega$ \BookRecordIndex{T0372}驯服，
则未必如此。不过，只要 $\omega$ 闭，在两种情形下 \BookRecordIndex{T0003}$J$-全纯曲线的能量都是拓扑不变量。
对带边黎曼曲面也有同样结论，只要所考虑的 \BookRecordIndex{T0003}$J$-全纯曲线 $u:\Sigma\to M$ 把边界映到
给定\BookRecordIndex{T0219}拉格朗日子流形 $L\subset M$ 中。这一观察在 \BookRecordIndex{T0003}$J$-全纯曲线的格罗莫夫紧致性定理
（第 4、5 章）中起关键作用。一般殆复流形中的 \BookRecordIndex{T0003}$J$-全纯曲线，与\BookRecordIndex{T0088}殆复结构被\BookRecordIndex{T0363}辛形式\BookRecordIndex{T0372}驯服的
流形中的 \BookRecordIndex{T0003}$J$-全纯曲线，其关键差别就在这里。本章关于 \BookRecordIndex{T0003}$J$-全纯曲线局部行为的其余结果，
都只涉及一般殆复流形。

\begin{exercise}\label{ex:2-2-2}
在引理 2.2.1 的记号下，证明：若 $J$ 被 $\omega$ \BookRecordIndex{T0372}驯服且 $u$ 是 $J$-全纯的，则
\begin{equation}\label{eq:2-2-4}
  \frac12|du|_J^2d\operatorname{vol}_\Sigma
  =\omega(\partial_su,J\partial_su)\,ds\wedge dt.
\end{equation}
\end{exercise}

\begin{exercise}\label{ex:2-2-3}
设 $L:\mathbb C\to\mathbb R^{2n}$ 是任意实线性映射。证明表达式
\[
  |L|:=|\zeta|^{-1}\sqrt{|L(\zeta)|^2+|L(i\zeta)|^2}
\]
与 $\zeta\in\mathbb C\setminus\{0\}$ 的选择无关，其中 $|\cdot|$ 表示
$\mathbb R^{2n}$ 上的欧氏范数。若同一化 $\mathbb R^{2n}=\mathbb C^n$ 且 $L$ 复线性，
则 $|L|=\sqrt2|L(1)|$；证明 $|L(1)|$ 是 $L$ 的算子范数，并证明 $|L|$ 是在标准正交基下
表示 $L$ 的实矩阵各项平方和的平方根。若 $\phi:\mathbb C\to\mathbb C$ 是全纯函数，
证明 $|d\phi(z)|=\sqrt2|\phi'(z)|$，其中 $\phi'(z)\in\mathbb C$ 表示其复导数。
\end{exercise}
